\section{Background and Problem}
\label{sec:background}

\paragraph{Synchronous data parallelism over a WAN.}
In DDP~\citep{fsdp} every step ends with an all-reduce of the full gradient. Gloo's ring all-reduce over
$n$ ranks runs $2(n-1)$ synchronous rounds; in each round every rank sends $S/n$ bytes to its successor
in rank order. A round finishes when its slowest ring edge finishes, so step time depends on the
\emph{worst} inter-site edge the ring order creates. That edge's cost depends on the path routing picks
for it, and on whatever else (background traffic, checkpoint uploads) shares that path. Sharded variants
(FSDP2, HSDP) trade all-reduce for all-gather/reduce-scatter but face the same coupling.
Placement, routing and checkpoint traffic therefore interact, even though they are usually owned by
different teams.

\paragraph{Traffic engineering without the job.}
WAN TE systems~\citep{b4,swan} allocate measured demand to paths to keep utilisation low and fair. A TE
optimiser that minimises maximum link utilisation can move a latency-sensitive collective onto a
higher-capacity, higher-latency path, which is good for utilisation and bad for a synchronous ring. It also
does not know that a burst of checkpoint traffic will arrive every $K$ steps, or that the job could move.

\paragraph{Checkpointing.}
CheckFreq~\citep{checkfreq} splits checkpointing into a snapshot and an asynchronous persist phase.
Check-N-Run~\citep{checknrun} uses incremental, quantized checkpoints, and Gemini~\citep{gemini}
checkpoints to remote CPU memory and schedules checkpoint traffic around training traffic. In a
geo-distributed job, the checkpoint destination $s$ is itself a WAN decision constrained by residency
policy: here, checkpoints of models trained on Canadian confidential data must stay in Toronto or
Montr\'eal.

\paragraph{Problem statement.}
Given a job (ranks, gradient bytes $S$, compute time $C$, checkpoint size $B$ every $K$ steps) and a
time-varying network state (link capacity, delay, up/down, background load), choose
$\pi$ (rank$\to$site in ring order), $r$ ((src site, dst site, traffic class)$\to$plane) and
$s \in \{\text{tor},\text{mtl}\}$ to minimise expected time per step, including amortised checkpoint
stalls. Re-planning is allowed: route changes are cheap, and re-placement costs a drain, restart and
restore. Our primary scope is $(\pi, r, s)$. The general objective in the project brief also includes
inference, media and storage-cost terms; we evaluate media separately (\S\ref{sec:ext}) and do not
optimise those terms jointly.
